\chapter{Hysterectomy}

\section*{Introduction \& Clinical Indications}
Hysterectomy is a significant surgical intervention involving the complete or partial excision of the uterus, a vital reproductive organ located in the female pelvis. This procedure is primarily indicated for various gynecological conditions, including both benign and malignant disorders, particularly when conservative management has failed or is not appropriate. The surgery can be categorized into total hysterectomy, which involves the removal of the uterus and cervix, and subtotal (or supracervical) hysterectomy, where the cervix is preserved. In some cases, additional structures such as the ovaries and fallopian tubes may also be removed, a procedure known as oophorectomy and salpingectomy, respectively. Hysterectomy plays a crucial role in managing conditions like uterine fibroids, endometriosis, abnormal uterine bleeding, and gynecological cancers, significantly impacting a woman's reproductive health and hormonal balance. The surgical approach can vary, including abdominal, vaginal, laparoscopic, or robotic-assisted techniques, each with unique benefits and considerations regarding recovery and potential complications.

\begin{itemize}
  \item Uterus removal surgery
  \item Total abdominal hysterectomy (TAH)
  \item Vaginal hysterectomy (VH)
  \item Laparoscopic hysterectomy (LH)
  \item Robotic-assisted hysterectomy (RAH)
  \item Supracervical hysterectomy (subtotal hysterectomy)
\end{itemize}

The primary surgical principle of hysterectomy is the definitive excision of the uterus to eliminate pathological tissue, alleviate severe symptoms, or treat malignancy. The operative concept involves a systematic dissection of the uterus from its surrounding anatomical attachments and vascular supply while preserving adjacent vital structures. The procedure typically begins with the division and ligation of the round ligaments, followed by the infundibulopelvic ligaments if oophorectomy is performed, or the ovarian ligaments and fallopian tubes. A critical step is the careful identification, clamping, division, and ligation or coagulation of the uterine arteries and veins within the cardinal ligaments, ensuring complete hemostasis and devascularization of the uterus.

The rationale for performing a hysterectomy arises from the failure of conservative medical or less invasive surgical treatments to adequately manage severe gynecological conditions. For benign diseases such as symptomatic uterine fibroids, adenomyosis, or intractable abnormal uterine bleeding, hysterectomy offers a permanent solution by removing the source of pathology. In cases of gynecological malignancy, such as endometrial or cervical cancer, the rationale is oncological eradication, aiming to remove all cancerous tissue and achieve clear surgical margins, often necessitating a more extensive resection, including parametrial tissue and pelvic lymph nodes. The ultimate technical goal is to achieve complete removal of the diseased uterus with minimal morbidity, preserving pelvic floor integrity where possible, and optimizing the patient's long-term quality of life.

The history of hysterectomy is marked by significant milestones that reflect advancements in surgical techniques and patient safety. The first recorded successful abdominal hysterectomy was performed by Konrad Langenbeck in Germany in 1813; however, the patient died shortly after due to complications. Charles Clay, an English surgeon, is credited with the first successful subtotal hysterectomy for a large fibroid uterus in 1843, with the patient surviving the procedure. The introduction of antiseptic techniques by Joseph Lister in the mid-19th century and advancements in anesthesia significantly improved surgical outcomes, making major abdominal surgery more feasible.

The early 20th century saw further refinement of techniques, with Ernst Wertheim's radical hysterectomy for cervical cancer, first performed in 1900, becoming the gold standard for decades. This technique emphasized the importance of parametrial and lymph node dissection for oncological cure. The latter half of the 20th century marked a paradigm shift with the advent of minimally invasive surgery. Harry Reich performed the first laparoscopic hysterectomy in 1988, demonstrating that the uterus could be safely removed through small abdominal incisions. This innovation paved the way for the widespread adoption of laparoscopic and robotic-assisted techniques, which offer significant benefits such as reduced postoperative pain, shorter hospital stays, and faster recovery compared to traditional open abdominal surgery, fundamentally transforming the approach to hysterectomy.

Hysterectomy is indicated for a variety of clinical conditions, including:

\begin{itemize}
  \item Symptomatic uterine fibroids (leiomyomas) causing heavy menstrual bleeding, chronic pelvic pain, or pressure symptoms on adjacent organs, refractory to conservative management.
  \item Severe endometriosis or adenomyosis causing chronic pelvic pain or abnormal uterine bleeding, when other therapies have failed and fertility preservation is not desired.
  \item Persistent, severe abnormal uterine bleeding (AUB) unresponsive to medical therapy, particularly in perimenopausal or postmenopausal women where malignancy is a concern.
  \item Severe uterine prolapse causing significant discomfort or urinary/bowel dysfunction, necessitating pelvic floor repair.
  \item Gynecological cancers, including endometrial and cervical cancer, as a primary surgical intervention or part of a comprehensive oncological management plan.
  \item Chronic pelvic pain definitively attributed to the uterus when all other diagnostic and therapeutic avenues have been exhausted.
  \item Life-threatening obstetric complications, such as intractable postpartum hemorrhage unresponsive to uterotonics or placenta accreta spectrum disorders.
  \item Precancerous conditions of the cervix or endometrium, such as high-grade cervical intraepithelial neoplasia (CIN 3) or atypical endometrial hyperplasia, when conservative treatments are not feasible.
\end{itemize}

Hysterectomy serves a dual clinical positioning, acting as both a primary and secondary treatment depending on the underlying pathology and clinical urgency. In cases of confirmed or highly suspected gynecological malignancies, such as endometrial cancer or advanced cervical cancer, hysterectomy is often the primary surgical intervention aimed at disease eradication. In life-threatening obstetric emergencies, such as intractable postpartum hemorrhage, hysterectomy is performed urgently to control bleeding and prevent maternal mortality.

Conversely, for benign indications like symptomatic uterine fibroids or endometriosis, hysterectomy is typically considered a secondary or definitive treatment. It is usually reserved for patients who have failed multiple conservative medical therapies or less invasive surgical interventions. The decision to proceed with hysterectomy in these cases is often a collaborative decision between the patient and surgeon, weighing the benefits of symptom resolution against the irreversibility and potential risks of the surgery.

\section*{Relevant Surgical Anatomy}
A comprehensive understanding of the surgical anatomy of the female pelvis is essential for performing a hysterectomy safely and effectively. The uterus is a hollow, muscular organ that typically measures approximately 7-8 cm in length and 5 cm in width in nulliparous women. It consists of three main parts: the fundus (the upper dome), the corpus (the body), and the cervix, which connects to the vagina. The uterus is supported by several ligaments, including the broad ligaments, round ligaments, cardinal ligaments, and uterosacral ligaments. 

The broad ligaments are peritoneal folds that extend from the lateral aspects of the uterus to the pelvic sidewalls, containing the fallopian tubes and ovarian ligaments. The round ligaments extend from the uterine fundus to the labia majora, providing anterior support. The cardinal ligaments, also known as Mackenrodt's ligaments, are critical for cervical support and contain the uterine artery and vein. The uterosacral ligaments extend from the cervix to the sacrum, providing posterior support.

The arterial supply to the uterus is primarily via the uterine arteries, branches of the internal iliac arteries, which run medially within the cardinal ligaments and cross anterior to the ureters. The ovarian arteries, originating from the aorta, also contribute to the blood supply. Venous drainage parallels the arterial supply, with the uterine veins draining into the internal iliac veins. Lymphatic drainage is complex, primarily involving the external iliac, internal iliac, obturator, and sacral lymph nodes. The bladder is located anterior to the uterus, separated by the vesicouterine pouch, while the rectum lies posterior, separated by the rectouterine pouch (pouch of Douglas). The close proximity of the ureters, bladder, and bowel necessitates meticulous dissection to prevent iatrogenic injury.

\section*{Preoperative Workup \& Preparation}
A thorough preoperative workup and preparation are essential to ensure patient safety, optimize surgical outcomes, and minimize the risk of complications. This comprehensive process typically involves a detailed medical and surgical history, physical examination, and targeted investigations.

Key components of preoperative preparation include:

\begin{itemize}
  \item \textbf{Medical and Cardiac Evaluation:} A comprehensive history and physical examination, including a pelvic exam and Pap smear (if indicated), are performed. Pre-existing medical conditions (e.g., cardiac disease, diabetes, respiratory issues, hypertension) are thoroughly assessed, optimized, and managed by the primary care physician or relevant specialists. Anesthesia consultation is mandatory to assess fitness for general anesthesia.

  \item \textbf{Laboratory Tests:} Routine blood tests include a complete blood count (CBC) to assess for anemia, blood type and screen for potential transfusion, coagulation profile (PT/INR, PTT) to evaluate bleeding risk, and kidney and liver function tests. Urinalysis and urine culture are often performed to rule out asymptomatic urinary tract infection, which should be treated preoperatively.

  \item \textbf{Imaging Studies:} Depending on the indication, imaging such as pelvic ultrasound, MRI, or CT scan may be performed to assess uterine size, fibroid location, extent of endometriosis, or staging of malignancy.

  \item \textbf{Medication Review and Adjustment:} Patients are instructed to stop certain medications, such as aspirin, NSAIDs, and blood thinners (e.g., warfarin, direct oral anticoagulants), typically 7-10 days prior to surgery to minimize bleeding risk. Other medications, such as diabetes medications, may need adjustment. Herbal supplements should also be reviewed and often discontinued.

  \item \textbf{Bowel Preparation:} For some approaches (e.g., vaginal, radical hysterectomy) or in cases of suspected extensive adhesions or bowel involvement, a mechanical bowel preparation combined with oral antibiotics may be prescribed to reduce the risk of infection if bowel injury occurs.

  \item \textbf{NPO Status:} Patients are instructed to fast (nil per os - NPO) for at least 6-8 hours before surgery (no food after midnight, clear liquids up to 2 hours before) to prevent aspiration during anesthesia.

  \item \textbf{Antibiotic Prophylaxis:} Intravenous broad-spectrum antibiotics (e.g., a cephalosporin) are typically administered within 60 minutes prior to the skin incision to reduce the risk of surgical site infection.

  \item \textbf{Informed Consent:} A detailed discussion of the procedure, its specific steps, potential risks, expected benefits, and available alternatives is conducted, followed by obtaining fully informed consent from the patient.
\end{itemize}

\textbf{Situations requiring diagnostic clarification, fertility counseling, optimization, or an alternative plan:}

\begin{itemize}
  \item \textbf{Desire for Future Fertility:} For benign indications, if a woman wishes to preserve her ability to conceive, hysterectomy is absolutely contraindicated. Alternative fertility-sparing treatments (e.g., myomectomy, endometrial ablation, medical management) must be thoroughly explored and offered.

  \item \textbf{Severe Uncontrolled Medical Instability:} Patients with acute, severe, and uncorrected medical conditions (e.g., uncontrolled sepsis, severe cardiac decompensation, acute respiratory failure, recent myocardial infarction or stroke) that render general anesthesia or major surgery prohibitively risky. These conditions must be stabilized before elective surgery.

  \item \textbf{Active, Untreated Pelvic Infection:} An active, untreated pelvic infection (e.g., acute pelvic inflammatory disease, tubo-ovarian abscess) is generally an absolute contraindication due to the high risk of spreading infection, bacteremia, and sepsis, unless the hysterectomy is performed as a life-saving measure to remove the source of infection (e.g., ruptured abscess).
\end{itemize}

\textbf{Relative Contraindications \& Precautions:}

\begin{itemize}
  \item \textbf{Morbid Obesity:} While not an absolute contraindication, morbid obesity significantly increases surgical complexity, operative time, and risks of complications (e.g., wound infection, deep vein thrombosis, pulmonary embolism, anesthesia complications). Minimally invasive approaches may be preferred, but technical challenges can still arise.

  \item \textbf{Extensive Pelvic Adhesions:} A history of previous abdominal or pelvic surgeries, severe endometriosis, or pelvic inflammatory disease can lead to dense adhesions, increasing the risk of injury to adjacent organs (bowel, bladder, ureters) and making the surgery more challenging and prolonged. Careful preoperative imaging and patient counseling are crucial.

  \item \textbf{Coagulopathy or Bleeding Disorders:} Uncorrected bleeding disorders (e.g., hemophilia, von Willebrand disease) or ongoing anticoagulant therapy require careful management, including hematology consultation and potential temporary cessation or reversal of anticoagulation, to minimize intraoperative and postoperative hemorrhage risk.

  \item \textbf{Extremely Large Uterine Size:} An extremely large uterus (e.g., exceeding 20 weeks' gestation size or weighing over 1000 grams) can make vaginal or laparoscopic approaches technically difficult or impossible, often necessitating an open abdominal approach.

  \item \textbf{Severe Anemia:} While not a contraindication to surgery itself, severe anemia should ideally be corrected preoperatively with iron supplementation or blood transfusion to minimize the need for intraoperative blood transfusions and improve patient tolerance to blood loss.

  \item \textbf{Lack of Surgical Expertise for Chosen Approach:} The chosen surgical approach (abdominal, vaginal, laparoscopic, robotic) should always match the surgeon's experience, skill set, and available institutional resources, especially for complex cases or advanced minimally invasive techniques. Referral to a more experienced surgeon may be indicated.
\end{itemize}

\section*{Surgical Technique \& Procedure}
Hysterectomy is performed under general anesthesia, with patient positioning varying by approach: dorsal lithotomy for vaginal, laparoscopic, or robotic approaches, or supine for an open abdominal approach. Prophylactic antibiotics are administered intravenously prior to incision.

The general steps for a total hysterectomy (removal of uterus and cervix) are as follows:

\begin{enumerate}
  \item \textbf{Incision and Access:}
    \begin{itemize}
      \item \textit{Abdominal:} A transverse Pfannenstiel incision (curved, just above the pubic hairline) or a vertical midline incision (from umbilicus to pubic symphysis) is made to enter the peritoneal cavity.
      \item \textit{Laparoscopic/Robotic:} Several small incisions (5-12 mm ports) are made in the abdomen. Carbon dioxide gas is insufflated to create pneumoperitoneum, and a camera port is inserted for visualization.
      \item \textit{Vaginal:} Incisions are made within the vagina, typically circumferentially around the cervix.
    \end{itemize}
  
  \item \textbf{Identification and Mobilization:} The uterus and adnexa are identified. For abdominal and minimally invasive approaches, the bowel is packed away to optimize visualization.

  \item \textbf{Division of Round and Ovarian/Infundibulopelvic Ligaments:} The round ligaments are identified, clamped, cut, and ligated or sealed with energy devices. If oophorectomy is planned, the infundibulopelvic ligaments (containing the ovarian artery and vein) are similarly divided. If ovaries are preserved, the ovarian ligaments and fallopian tubes are divided close to the uterus.

  \item \textbf{Creation of Bladder Flap:} The peritoneum overlying the lower uterine segment is incised, and the bladder is carefully dissected inferiorly off the anterior uterine wall and cervix, creating a vesicouterine space to protect the bladder and expose the uterine vessels.

  \item \textbf{Ligation of Uterine Arteries and Veins:} The uterine arteries and veins, which run within the cardinal ligaments, are meticulously identified. These vessels are clamped, divided, and ligated or sealed with energy devices, ensuring complete hemostasis. Care is taken to avoid injury to the ureters, which pass close to the uterine arteries.

  \item \textbf{Division of Uterosacral and Cardinal Ligaments:} The uterosacral and cardinal ligaments, which provide significant support to the cervix, are systematically clamped, cut, and ligated or sealed, further mobilizing the uterus.

  \item \textbf{Cervical Transection and Vaginal Cuff Closure (Total Hysterectomy):} The cervix is carefully separated from the upper vagina. The vaginal cuff (the top of the vagina) is then closed with sutures, often in a continuous or interrupted fashion, to prevent prolapse and provide support. In supracervical hysterectomy, the cervix is left intact, and the uterine body is transected above the cervix.

  \item \textbf{Uterus Removal and Hemostasis:} The uterus (and cervix, if total) is removed. In abdominal cases, it is removed through the incision. In vaginal cases, it is delivered through the vagina. In laparoscopic/robotic cases, it may be removed vaginally, morcellated (cut into smaller pieces) and removed through a port site, or removed through an enlarged port site. The surgical field is meticulously inspected for any bleeding, and hemostasis is ensured.

  \item \textbf{Closure:} The abdominal incisions or vaginal cuff are closed in layers. Drains are rarely used but may be placed in complex cases or when significant fluid collection is anticipated. The skin is closed with sutures or staples.
\end{enumerate}

\section*{Complications \& Risks}
As with any major surgical procedure, hysterectomy carries potential risks, which can be broadly categorized into general surgical complications and procedure-specific risks.

\begin{itemize}
  \item \textbf{General Surgical Complications:}
    \begin{itemize}
      \item \textbf{Bleeding/Hemorrhage:} Intraoperative or postoperative hemorrhage, which may necessitate blood transfusion, reoperation, or, in rare severe cases, lead to hypovolemic shock.
      \item \textbf{Infection:} Surgical site infection (e.g., wound infection, vaginal cuff cellulitis), pelvic infection (e.g., abscess), or urinary tract infection (UTI) due to catheterization.
      \item \textbf{Anesthesia Complications:} Adverse reactions to general anesthesia, including respiratory depression, cardiac arrhythmias, aspiration pneumonitis, nausea, vomiting, or allergic reactions.
      \item \textbf{Deep Vein Thrombosis (DVT) and Pulmonary Embolism (PE):} Formation of blood clots in the deep veins of the legs, which can dislodge and travel to the lungs, a potentially life-threatening complication. Prophylaxis (e.g., sequential compression devices, anticoagulants) is routinely used.
      \item \textbf{Pain:} Postoperative pain, which is managed with analgesics but can occasionally be severe or persistent.
    \end{itemize}

  \item \textbf{Procedure-Specific Risks:}
    \begin{itemize}
      \item \textbf{Injury to Adjacent Organs:} Damage to the bladder, ureters (tubes connecting kidneys to bladder), or bowel during dissection. Such injuries may require immediate surgical repair, potentially leading to prolonged recovery, fistula formation (e.g., vesicovaginal, rectovaginal), or peritonitis.
      \item \textbf{Vaginal Cuff Dehiscence or Prolapse:} The surgical incision at the top of the vagina (vaginal cuff) can separate (dehiscence), requiring reoperation. In the long term, the vaginal vault can prolapse due to the loss of uterine support, necessitating further reconstructive surgery.
      \item \textbf{Nerve Injury:} Damage to nerves in the pelvis or legs (e.g., femoral nerve, obturator nerve, sciatic nerve) due to prolonged positioning, retraction, or direct surgical manipulation, leading to numbness, weakness, or chronic neuropathic pain.
      \item \textbf{Formation of Adhesions:} Scar tissue can form internally between organs, potentially causing chronic pelvic pain, bowel obstruction, or infertility in the future.
      \item \textbf{Early Menopause (if ovaries removed):} If both ovaries are removed (bilateral oophorectomy) in a premenopausal woman, surgical menopause occurs abruptly, leading to symptoms like hot flashes, vaginal dryness, mood swings, and increased long-term risk of osteoporosis and cardiovascular disease.
      \item \textbf{Sexual Dysfunction:} While many women report improved sexual function after symptom resolution, some may experience changes in sexual sensation, libido, or discomfort during intercourse.
      \item \textbf{Psychological Impact:} Emotional distress, depression, anxiety, or feelings of loss, particularly in younger women or those who desired future fertility.
      \item \textbf{Chronic Pain:} Persistent pelvic pain despite hysterectomy, especially if the original cause was complex, multifactorial, or not solely uterine in origin.
      \item \textbf{Urinary Dysfunction:} New onset or worsening of urinary incontinence (stress or urge) or urinary urgency due to changes in pelvic anatomy or nerve function.
      \item \textbf{Fistula Formation:} Abnormal connections between the vagina and bladder (vesicovaginal fistula) or rectum (rectovaginal fistula) can occur, leading to continuous leakage of urine or stool.
    \end{itemize}
\end{itemize}

\section*{Postoperative Care \& Recovery}
Following hysterectomy, patients are transferred to the Post-Anesthesia Care Unit (PACU) for close monitoring of vital signs, pain levels, and recovery from anesthesia. Pain management is initiated, often with intravenous opioids, and gradually transitioned to oral pain relievers. Early ambulation is strongly encouraged within hours of surgery to prevent complications such as deep vein thrombosis, stimulate bowel function, and aid in pulmonary recovery.

During the first 24-48 hours, patients typically remain in the hospital. The length of stay varies by surgical approach: laparoscopic and vaginal hysterectomy patients often have shorter hospital stays (1-2 days) compared to abdominal hysterectomy patients (2-4 days). A urinary catheter may be in place temporarily. Diet is advanced from clear liquids to solid foods as tolerated, usually once bowel sounds return. Patients are closely monitored for signs of complications such as excessive vaginal bleeding, fever, severe abdominal pain, or signs of infection. Before discharge, patients receive detailed instructions on wound care, activity restrictions (e.g., no heavy lifting, no strenuous exercise, no vaginal intercourse), pain medication regimen, and warning signs to watch for. Full recovery at home typically takes 4-6 weeks for minimally invasive approaches and 6-8 weeks for abdominal hysterectomy, during which internal tissues heal and energy levels gradually return. Emotional support and counseling may be beneficial for some women adjusting to the physical and emotional changes.

During the hysterectomy, the surgeon meticulously documents intraoperative findings, including the size and appearance of the uterus, presence and characteristics of fibroids, extent of endometriosis or adenomyosis, presence of adhesions, and any suspicious lesions on the ovaries or fallopian tubes. The condition of adjacent organs like the bladder and bowel is also noted. Following removal, the resected uterine specimen (and any adnexa) is sent to the pathology laboratory for gross and microscopic examination. The pathology report is the definitive "result" of the surgery. It details the macroscopic description of the uterus (size, weight, presence of masses) and, crucially, the microscopic findings. This report confirms the preoperative diagnosis (e.g., leiomyomas, adenomyosis, endometrial hyperplasia) or provides the definitive diagnosis for malignancy. For cancer cases, the pathology report specifies the exact type of cancer, its histological grade, depth of invasion into the myometrium, involvement of the cervix or other structures, and the status of surgical margins (whether all cancerous tissue was removed with a clear border of healthy tissue). If lymph nodes were removed, their status (positive or negative for cancer cells) is also reported, which is vital for cancer staging and guiding any subsequent adjuvant therapy.

A normal and successful postoperative course following a hysterectomy is characterized by a progressive resolution of symptoms and a gradual return to normal activities without significant complications. Patients can expect initial incisional pain or vaginal discomfort, which should be well-controlled with prescribed oral analgesics and steadily improve over the first few weeks. Light vaginal spotting or discharge is common for several weeks as the vaginal cuff heals. Bowel function typically returns within 1-3 days, and patients should experience a gradual increase in energy levels. By 2-3 weeks, most patients are able to resume light daily activities, and by 4-6 weeks (or 6-8 weeks for abdominal approaches), they are generally cleared to resume all normal activities, including exercise and sexual intercourse, once the vaginal cuff is fully healed. The primary symptoms that prompted the surgery, such as heavy bleeding or pelvic pain, should be completely resolved. For women who had their ovaries removed, menopausal symptoms will begin, which can be managed with hormone replacement therapy if appropriate.

Postoperative care and follow-up are crucial for monitoring recovery, addressing any concerns, and ensuring optimal long-term health.

\begin{itemize}
  \item \textbf{First Postoperative Visit:} Typically scheduled 2-3 weeks after surgery to check wound healing, remove any external sutures or staples, and discuss initial recovery progress and pain management.

  \item \textbf{Second Postoperative Visit:} Around 6 weeks post-surgery, a more comprehensive check-up is performed, including a pelvic exam to assess vaginal cuff healing and ensure there are no signs of infection or dehiscence. Activity restrictions are usually lifted at this point.

  \item \textbf{Pathology Discussion:} A detailed discussion of the pathology report findings and their implications for future health or further treatment, especially in cases of malignancy, is conducted.

  \item \textbf{Hormone Replacement Therapy (HRT):} If both ovaries were removed in premenopausal women, discussion and initiation of HRT may occur to manage menopausal symptoms and mitigate long-term health risks such as osteoporosis and cardiovascular disease.

  \item \textbf{Cervical Cancer Screening:} If the cervix was preserved (supracervical hysterectomy), regular Pap smears are still necessary according to national guidelines. If the cervix was removed (total hysterectomy), routine Pap smears are generally no longer needed, but vaginal cuff cytology may be recommended in specific cases (e.g., history of high-grade cervical dysplasia or cancer).

  \item \textbf{Pelvic Floor Physical Therapy:} May be recommended for women experiencing new or persistent urinary incontinence, pelvic organ prolapse symptoms, or other pelvic floor dysfunction following surgery.

  \item \textbf{Warning Signs:} Patients are advised to contact their doctor immediately for signs of complications such as fever (over 100.4$^{\circ}$F or 38$^{\circ}$C), severe or worsening abdominal pain, heavy vaginal bleeding (soaking more than one pad per hour), foul-smelling vaginal discharge, redness, swelling, or pus at the incision site, or symptoms of DVT (leg pain, swelling).
\end{itemize}

\section*{Outcomes \& Prognosis}
Hysterectomy remains one of the most common major surgical procedures performed on women globally. In the United States, approximately 400,000 to 500,000 hysterectomies are performed annually, although the incidence has seen a gradual decline over the past few decades due to the increasing availability and adoption of less invasive alternatives for benign conditions. The procedure is most frequently performed in women between the ages of 40 and 50, with the average age being around 45 years, reflecting the peak incidence of conditions like symptomatic fibroids and abnormal uterine bleeding. The leading indications for hysterectomy are uterine fibroids (leiomyomas), accounting for approximately 30-40\% of cases, followed by abnormal uterine bleeding, endometriosis, adenomyosis, and uterine prolapse. Gynecological cancers (endometrial, cervical, ovarian) account for a smaller but significant proportion, typically 10-15\% of all hysterectomies. While overall mortality rates are very low (typically less than 0.1\% in elective cases), morbidity, including complications such as infection, hemorrhage, and organ injury, can range from 5\% to 10\%, depending on the surgical approach, patient comorbidities, and surgeon experience. Minimally invasive approaches (laparoscopic, robotic, vaginal) generally demonstrate lower morbidity rates, shorter hospital stays, and faster recovery compared to open abdominal hysterectomy.

Hysterectomy can definitively treat selected uterine conditions, but symptom response varies by indication and extrauterine disease. Counseling should address fertility loss, cervical and ovarian management, pelvic-floor effects, sexual function, and possible earlier menopause if ovarian function is affected; no universal symptom-resolution percentage applies. For cancer, prognosis depends on stage, histology, margins, nodes, and multimodal treatment.

\section*{References}
\begin{enumerate}
  \item MedlinePlus. Hysterectomy. U.S. National Library of Medicine; 2024. Available from: \url{https://medlineplus.gov/hysterectomy.html}
  
  \item Hoffman BL, Schorge JO, Bradshaw KD, Halvorson LM, Schaffer JI, Corton LD, editors. Williams Gynecology. 4th ed. New York, NY: McGraw-Hill Education; 2024.
  
  \item American College of Obstetricians and Gynecologists (ACOG). Hysterectomy. Patient Education Pamphlet; 2023.
  
  \item Sabiston DC, Townsend CM, Beauchamp RD, Evers BM, Mattox KL, editors. Sabiston Textbook of Surgery: The Biological Basis of Modern Surgical Practice. 21st ed. Philadelphia, PA: Elsevier; 2022.
\end{enumerate}

\clearpage
